% Generated draft: requires mathematical and source review.
\documentclass[11pt]{article}
\usepackage[margin=25mm]{geometry}
\usepackage{fontspec}
\setmainfont{Latin Modern Roman}
\usepackage{amsmath,amssymb,mathtools,unicode-math}
\setmathfont{Latin Modern Math}
\usepackage{enumitem,needspace,xurl,hyperref}
\setlength{\parindent}{0pt}
\setlength{\parskip}{5pt}
\setlength{\emergencystretch}{3em}
\newcommand{\N}{\mathbb{N}}
\newcommand{\Z}{\mathbb{Z}}
\newcommand{\Q}{\mathbb{Q}}
\newcommand{\R}{\mathbb{R}}
\newcommand{\C}{\mathbb{C}}
\newcommand{\F}{\mathbb{F}}
\newcommand{\A}{\mathbb{A}}
\newcommand{\PP}{\mathbb{P}}
\newcommand{\E}{\mathbb{E}}
\newcommand{\eps}{\varepsilon}
\newcommand{\dd}{\,\mathrm{d}}
\begin{document}
\section*{UnsolvedMath ID: 30002329}
\textbf{Source problem code:} catalogue entry
\subsection{Definitions and mathematical statement}
The supplied question concerns global Willmore minimization near the conformal class of the Clifford torus. Its precise admissible class and the meaning of ``constrained energy'' are not specified. The following conventional framework makes the intended comparison explicit, but its applicability to the source requires review.

Let \(T^2=\R^2/\Z^2\), and let \(S^3=\{(z_1,z_2)\in\C^2:|z_1|^2+|z_2|^2=1\}\) have its round metric. For a smooth immersion \(f:T^2\to S^3\), define
\[
\mathcal W(f)=\int_{T^2}(1+H_f^2)\,\dd A_f,
\]
where \(H_f\) is the average of the two principal curvatures and \(\dd A_f\) is the induced area measure. Let \([f]\) denote the conformal class of the induced metric, viewed in the moduli space \(\mathcal M_1\) of complex tori. Set
\[
f_{\mathrm{Cl}}(x,y)=\frac{1}{\sqrt2}\bigl(e^{2\pi i x},e^{2\pi i y}\bigr),
\qquad c_{\mathrm{Cl}}=[f_{\mathrm{Cl}}].
\]
For \(c\in\mathcal M_1\), put
\[
\mathcal I(c)=\{f:T^2\to S^3: f\text{ is a smooth immersion and }[f]=c\}.
\]
A constrained Willmore immersion is a critical point of \(\mathcal W\) under variations preserving its conformal class. A Hopf torus is obtained by lifting a closed regular curve through the Hopf fibration \(S^3\to\mathbb{CP}^1\).

Let \(\mathcal K(c)\) denote the source-described constrained Willmore Hopf tori of conformal class \(c\) that ``minimize the constrained energy.'' This designation does not yet define a precise candidate set: the supplied record gives neither that energy functional nor the class over which it is minimized.

Subject to these unresolved choices, the intended local-to-global question has the form: does there exist an open neighborhood \(U\subseteq\mathcal M_1\) of \(c_{\mathrm{Cl}}\) such that
\[
\forall c\in U\quad\forall h\in\mathcal K(c)\quad
\mathcal W(h)=\inf_{f\in\mathcal I(c)}\mathcal W(f)?
\]
Equivalently, must every such candidate have Willmore energy no greater than that of every admissible immersion in its conformal class? This formulation does not assert that \(\mathcal K(c)\) is nonempty for every \(c\in U\). The source does not resolve whether such existence is part of the intended assertion, whether every constrained-energy minimizer is intended, or whether a specified candidate family is intended. The ambient space, regularity, immersion versus embedding requirements, treatment of covers, and marked versus unmarked conformal classes also require source confirmation. Thus the displayed formula is a conditional formulation of the comparison, not a completed resolution of the missing hypotheses.
\subsection{Short English statement}
The question asks whether constrained-energy-minimizing Hopf tori with conformal classes sufficiently close to the Clifford torus also minimize Willmore energy among all admissible tori in the same conformal class. The constrained energy and its competitors are unspecified, so the precise candidate family and global comparison class require review.
\subsection{Sources}
\begin{itemize}
\item Supplied local UnsolvedMath registry record, permanent ID \(30002329\), title ``Global Willmore Minimizers Near the Clifford Torus,'' collection ``Oberwolfach Reports Open Problems,'' exact source problem code \texttt{catalogue entry}. The proposed year is supplied as 2013; no individual proposer is named.
\item Supplied report citation: \emph{Progress in Surface Theory (2014)}, Oberwolfach Reports, pp.\ 1253--1312, DOI \texttt{10.4171/owr/2013/21}, \url{https://doi.org/10.4171/owr/2013/21}. The literature-source list also labels this reference ``Oberwolfach Report 21/2013.'' The report-year labeling requires confirmation.
\item Supplied literature source: Heller--Ndiaye, \emph{First explicit constrained Willmore minimizers of non-rectangular conformal class}, \url{https://doi.org/10.1016/j.aim.2021.107804}.
\item Supplied literature source: \emph{Candidates for non-rectangular constrained Willmore minimizers}, \url{https://doi.org/10.1016/j.geomphys.2021.104221}. No authors are supplied for this entry.
\item Literature assessment dated exactly \texttt{2026-08-21}: the supplied curation labels the problem ``partially solved'' and states that existence and minimization are proved in several rectangular and non-rectangular conformal classes near the Clifford class, while the full local-to-global constrained-Willmore minimizer conjecture remains open. This is supplied dated curation derived from the OWR record, not an independently verified assessment; its time-sensitive status and cited claims require independent verification.
\item The registry labels the statement ``exact'' and reports that its mathematical content was checked against the official Oberwolfach report, with notation, context, and question form normalized. That verification claim is supplied provenance and has not been independently checked in this draft.
\end{itemize}
\end{document}
